\section{结构为何曾经有效：翻译与 Instruction Tuning 的自然实验}

上一章完成形式变换，本章问历史动机。Encoder-decoder 的额外结构并非任意复杂化：2017 年的核心任务是机器翻译，输入和目标来自不同语言；后来的学术 instruction-tuning 数据又常呈长输入、短答案。结构与数据分布吻合时，归纳偏置会显著提高有限预算下的效率。

\subsection{三项额外结构}

本节先看完整清单，再聚焦第一项。相对 decoder-only，encoder-decoder 假设输入与目标足够不同，目标应读取已经完全编码的输入，并且输入 token 之间适合全双向交互。三项假设分别对应参数、层间连接和 attention mask。

\lecturefigure{hyung-slide-052.jpg}{Encoder-decoder 相对 decoder-only 的三项额外结构}{Hyung Won Chung official deck, p.~52}

\lecturefigure{hyung-slide-053.jpg}{第一项假设：输入与目标不同到值得使用独立参数}{Hyung Won Chung official deck, p.~53}

\begin{importantbox}{判断一个 inductive bias 是否仍有效的三问}
\begin{enumerate}
  \item 它最初匹配了什么数据分布或任务约束？
  \item 当前数据、目标和交互方式是否仍满足这个约束？
  \item 若移除结构并增加规模，性能曲线在哪个预算点交叉？
\end{enumerate}
只回答第一问会把历史合理性误当成永久正确性。
\end{importantbox}

\subsection{机器翻译：独立参数曾经很自然}

2017 年 Transformer 面向 source-to-target translation。输入是源语言句子，目标是另一种语言的输出，二者词表、语序和生成角色不同。在只优化翻译的设定下，让 encoder 专注源语言、decoder 专注目标语言是一种高效先验。本节先承认这一历史合理性，再比较通用预训练改变了哪些条件，避免用今天的任务接口倒推旧设计“从一开始就错误”。

\lecturefigure{hyung-slide-054.jpg}{机器翻译中的 source/target 分布差异支持独立参数}{Hyung Won Chung official deck, p.~54}

\lecturefigure{hyung-slide-055.jpg}{通用语言模型要融合跨语言知识，独立参数的动机减弱}{Hyung Won Chung official deck, p.~55}

读第二页时要注意任务已经改变。现代 pretraining 不只学习英德翻译，而要在多语言文本中形成共享世界知识；同一事实可能用不同语言表达。若把输入与目标永久隔离，模型可能更难复用表示。另一方面，多语专门化仍可能有价值，因此这页提出的是 scale-era counterargument，不是独立参数绝对无用。

\teachervoice{00:57:51--00:59:09 的历史对比非常具体：当目标只是翻译时，给不同语言分开参数“很自然”；当目标变成一般知识学习时，讲者更希望把跨语言知识合并。Inductive bias 的评价随着 objective 改变。}

\subsection{FLAN：长输入、短目标让 encoder-decoder 意外占优}

Instruction tuning 提供了更细的自然实验。同一团队在 PaLM（decoder-only）与 T5（encoder-decoder）上做 fine-tuning；虽然大部分时间用于优化 PaLM，T5 的增益却更大。讲者没有把结果归因于架构普遍优越，而是回到数据长度分布寻找匹配关系。

\lecturefigure{hyung-slide-056.jpg}{Instruction finetuning 把大量学术任务统一成自然语言指令}{Hyung Won Chung official deck, p.~56}

\lecturefigure{hyung-slide-057.jpg}{FLAN 实验中 encoder-decoder 的 fine-tuning 增益更大}{Hyung Won Chung official deck, p.~57}

\begin{knowledgebox}{读图：比较的是增益，不是绝对模型胜负}
表格强调 instruction tuning 前后的 improvement。不能只看红框就得出 T5 在所有任务、所有规模上优于 PaLM；课堂结论是 encoder-decoder 的结构与该批学术数据特别匹配，因此 fine-tuning 信号利用率更高。
\end{knowledgebox}

\lecturefigure{hyung-slide-058.jpg}{学术数据常有长输入、短目标的长度分布}{Hyung Won Chung official deck, p.~58}

若输入长度随机变量为 $L_x$、目标长度为 $L_y$，该数据集大致满足
\[
\mathbb{E}[L_x]\gg \mathbb{E}[L_y].
\]
其中 $\mathbb{E}[L_x]$ 与 $\mathbb{E}[L_y]$ 分别是输入和目标平均长度。长题面、短标签使两侧分布明显不同，独立 encoder/decoder 参数就能专门化；但长文本生成和多轮聊天会缩小甚至反转这种不对称。

\teachervoice{00:59:46--01:01:10，Hyung Won 说团队花了 99\% 时间优化 PaLM，T5 却在几天内得到更大增益；他的解释是数据长度偏置与架构“意外对上了”。这段 teacher voice 把结果从品牌比较变成分布匹配分析。}

\begin{warningbox}{Academic benchmark 的可评分性会塑造训练分布}
长目标难以人工标注和自动评分，因此学术数据偏向长输入、短答案。可评分性不是用户价值的同义词：创作、代码、分析与对话往往需要长输出。若模型架构只对 benchmark 分布优化，部署任务改变后原有结构优势可能消失。
\end{warningbox}

\subsection{本章小结}

机器翻译和 FLAN 说明归纳偏置为什么能在特定时代产生真实收益：它把有限容量投向数据分布最需要的方向。但通用预训练、长生成和多轮对话改变了输入/目标关系，迫使研究者重新测试参数分离。历史不是为了嘲笑旧设计，而是为了识别设计成立的条件。

